% संपूर्ण मराठी ग्रंथातील प्रकरण 22: लिंडस्ट्रॉमचे प्रमेय
% हा संपादनयोग्य स्रोतखंड आहे. संकलनासाठी संपूर्ण स्रोतसंग्रहातील
% अक्षररूपे, पूर्वभाग, चिन्हव्याख्या आणि आकृती-स्रोत आवश्यक आहेत.
% मूळ: Open Logic Project; मजकूर आणि रूपांतर CC BY 4.0.
% यंत्रानुवाद, दुरुस्ती आणि तपासणी: OpenAI Codex —
% GPT-5.6 Sol आणि GPT-6 Sol, दोन्ही Ultra effort.
% दुरुस्ती-पुनरावलोकन: GPT-6.1 Sol, Ultra effort.
\chapter{लिंडस्ट्रॉमचे प्रमेय}\label{mod:lin::chap}








\section{परिचय}\label{mod:lin:int:sec}

काही आणखी निर्बंध गृहीत धरल्यास संहतता आणि अधोगामी
लोव्हेनहाइम--स्कोलेम प्रमेये ज्या महत्तम तर्कशास्त्रासाठी खरी ठरतात, ते
प्रथम-क्रम तर्कशास्त्र आहे, असे सांगणारे लिंडस्ट्रॉमचे लक्षणचित्रण या
प्रकरणात सिद्ध करणे हे आपले उद्दिष्ट आहे
(\ref{fol:com:com:thm:compactness} आणि
\ref{fol:com:dls:thm:downward-ls}). प्रथम, हे प्रमेय ज्या
तर्कशास्त्रांच्या सामान्य वर्गाला लागू होते, त्याचे अधिक व्यापक वर्णन
आपल्याला आवश्यक आहे. आपण स्वतःला \emph{संबंधात्मक} भाषांपुरते मर्यादित
ठेवू, म्हणजे ज्यांत केवळ विधेये आणि व्यक्ती सूचित करणारी स्थिर
चिन्हे असतात, पण फलने नसतात अशा भाषा.











\section{अमूर्त तर्कशास्त्रे}\label{mod:lin:alg:sec}

\begin{defn}
\emph{अमूर्त तर्कशास्त्र} म्हणजे जोडी $
\tuple{L, \models_L}$; येथे $L$ हे प्रत्येक भाषा~$\Lang{L}$
ला $L(\Lang{L})$ हा वाक्यांचा संच नेमून देणारे फलन आहे, आणि
$\models_L$ हा भाषा~$\Lang{L}$ साठीच्या संरचना आणि
$L(\Lang{L})$ च्या घटक यांमधील संबंध आहे. विशेषतः,
$\tuple{F, \models}$ हे साधारण प्रथम-क्रम तर्कशास्त्र आहे; म्हणजे $F$ हे
भाषा~$\Lang{L}$ ला $\Lang{L}$ मधील स्थिर चिन्हांपासून बनवलेल्या प्रथम-क्रम
वाक्यांचा संच त्या भाषेला नेमून देणारे फलन आहे, आणि $\models$ हा
प्रथम-क्रम तर्कशास्त्राचा पूर्ति-संबंध आहे.
\end{defn}

लक्षात घ्या की दिलेल्या भाषा साठी आपण प्रथम-क्रम तर्कशास्त्रातील
संरचनेचीच संकल्पना अजूनही वापरत आहोत; परंतु $\Lang{L}$ मधील
मूलभूत चिन्हांपासून वाक्ये नेहमीच्या रीतीने बनवली जातात, किंवा
$\models_L$ हा संबंध प्रथम-क्रम तर्कशास्त्राप्रमाणेच प्रत्यावर्ती रीतीने
परिभाषित केला जातो, असे आपण पूर्वगृहीत धरत नाही. म्हणून, ही पूर्णतः
सामान्य व्याख्या अशा प्रकरणांचाही समावेश करण्यासाठी आहे ज्यांत
$\tuple{L,\models_L}$ मधील वाक्ये मध्ये अनंत लांबीची संधी किंवा
विकल्प, $\lexists$ आणि $\lforall$ यांखेरीज इतर संख्यापक (उदा.,
``असे अनंत अनेक~$x$ आहेत की \dots''), किंवा कदाचित अनंत लांबीचे
संख्यापक-उपसर्ग असतात. $L(\Lang{L})$ मधील ``वाक्ये'' ही
प्रथम-क्रम तर्कशास्त्राची साधारण वाक्ये असतीलच असे नाही, हे
ठळक करण्यासाठी या प्रकरणात त्यांवर व्यापणारी चले म्हणून $\varphi_{E}$,
$\varphi_{F}$,~\dots वापरतो, आणि साधारण प्रथम-क्रम सूत्रे साठी $\varphi$,
$\psi$,~\dots राखून ठेवतो.

\begin{defn}
$\Mod(L){\varphi_{E}}$ ही संज्ञा $\Setabs{\Struct{M}}{\Struct{M}
  \models_L \varphi_{E}}$ हा वर्ग दर्शवू दे. भाषा स्पष्ट सांगणे आवश्यक
असेल, तर आपण $\Mod[L](L){\varphi_{E}}$ असे लिहितो. $\Lang{L}$ साठीच्या दोन
संरचना $\Struct{M}$ आणि $\Struct{N}$ या $\tuple{L, \models_L}$
मध्ये \emph{प्राथमिक सममूल्य} आहेत, आणि हे $\Struct{M} \elemequiv[L]
\Struct{N}$ असे लिहितो, जर $L(\Lang{L})$ मधील तीच वाक्ये
दोन्हींमध्ये सत्य असतील.
\end{defn}

\begin{defn}
भाषा $\Lang{L}$ साठीचे अमूर्त तर्कशास्त्र
$\tuple{L,\models_L}$ पुढील गुणधर्म पूर्ण करत असेल, तर ते \emph{सामान्य}
आहे:
\begin{enumerate}
\item (\emph{$L$-एकस्वनिकता}) भाषा $\Lang{L}$ आणि
  $\Lang{L'}$ यांच्यासाठी, $\Lang{L} \subseteq \Lang{L'}$ असेल, तर
  $L(\Lang{L}) \subseteq L(\Lang{L'})$.
\item (\emph{विस्तार गुणधर्म}) प्रत्येक $\varphi_{E} \in L(\Lang{L})$
  साठी $\Lang{L}$ चा असा \emph{सांत} उपसंच $\Lang{L'}$ असतो की
  $\Struct{M} \models_L \varphi_{E}$ हा संबंध $\Struct{M}$ च्या
  $\Lang{L'}$-न्यूनीकृत रचनेवरच अवलंबून असतो; म्हणजे $\Struct{M}$ आणि
  $\Struct{N}$ यांची $\Lang{L'}$ वरील न्यूनीकृत रचना एकच असेल, तर
  $\Struct{M} \models_L \varphi_{E}$ तेव्हा आणि केवळ तेव्हाच, जेव्हा
  $\Struct{N} \models_L \varphi_{E}$.
\item (\emph{समरूपता गुणधर्म}) $\Struct{M} \models_L \varphi_{E}$ आणि
  $\Struct{M} \simeq \Struct{N}$ असेल, तर $\Struct{N} \models_L \varphi_{E}$
  हेही खरे असते.
\item (\emph{पुनर्नामकरण गुणधर्म}) पुनर्नामकरणाखाली $\models_L$ संबंध
  टिकून राहतो: भाषा $\Lang{L}'$ ही $\Lang{L}$ मधील प्रत्येक
  चिन्ह $P$ च्या जागी त्याच स्थानसंख्येचे चिन्ह $P'$ आणि प्रत्येक स्थिर
  $c$ च्या जागी वेगळे स्थिर $c'$ ठेवून मिळत असेल, तर प्रत्येक
  संरचना~$\Struct{M}$ आणि वाक्य~$\varphi_{E}$ साठी,
  $\Struct{M} \models_L \varphi_{E}$ तेव्हा आणि केवळ तेव्हाच, जेव्हा
  $\Struct{M}' \models_L \varphi_{E}'$; येथे $\Struct{M}'$ ही $\Lang{L}'$-संरचना
  असून ती $\Lang{L}$ मधील मूळ रचनेशी अनुरूप आहे, आणि $\varphi_{E}' \in
  L(\Lang{L}')$.

\item (\emph{बूलीय गुणधर्म}) अमूर्त तर्कशास्त्र $\tuple{L,
  \models_L}$ बूलीय संयोजकांखाली पुढील अर्थाने संवृत असते: प्रत्येक
  $\varphi_{E} \in L(\Lang{L})$ साठी असा $\varphi_{F} \in L(\Lang{L})$ असतो की
  $\Struct{M} \models_L \varphi_{F}$ तेव्हा आणि केवळ तेव्हाच, जेव्हा
  $\Struct{M} \not\models_L \varphi_{E}$; आणि प्रत्येक $\varphi_{E}$ आणि $\varphi_{F}$ साठी असा
  $\varphi_{G}$ असतो की $\Mod(L){\varphi_{G}} = \Mod(L){\varphi_{E}} \cap
  \Mod(L){\varphi_{F}}$. आण्विक सूत्रे आणि इतर संयोजकांसाठीही अशाच अटी
  असतात.
\item (\emph{संख्यापक गुणधर्म}) $\Lang{L}$ मधील प्रत्येक स्थिर $c$
  आणि $\varphi_{E} \in L(\Lang{L})$ साठी असा $\varphi_{F} \in L(\Lang{L})$ असतो की
  \[  \Mod[L'](L){\varphi_{F}} = \Setabs{\Struct{M}}{\Expan{M}{a} \in
  \Mod[L](L){\varphi_{E}} \text{ काही } a \in \Domain{M}},
  \]
  येथे $\Lang{L'} = \Lang{L} \setminus \{c\}$ आणि
  $\Expan{M}{a}$ ही $c$ ला $a$ नेमून देणारी $\Struct{M}$ ची
  $\Lang{L}$ वरील विस्तारित रचना आहे.
\item (\emph{सापेक्षीकरण गुणधर्म}) $\varphi_{E} \in L(\Lang{L})$ हे
  वाक्य आणि $R$, $c_1$, \dots, $c_n$ ही $\Lang{L}$ मध्ये
  नसलेली चिन्हे दिली असता, असे वाक्य
  $\varphi_{F} \in L(\Lang{L} \cup \{R,c_1,\ldots,c_n\})$ असते की त्याला $\varphi_{E}$ चे
  $\Atom{R}{x, c_1, \dots c_n}$ पर्यंतचे \emph{सापेक्षीकरण} म्हणतात,
  आणि प्रत्येक संरचना~$\Struct{M}$ साठी:
  \[  \Expan{M}{X, b_1, \ldots, b_n} \models_L \varphi_{F} \text{ तेव्हा आणि
    केवळ तेव्हाच, जेव्हा } \Struct{N} \models_L \varphi_{E},
  \]
  येथे $\Struct{N}$ ही $\Struct{M}$ ची अशी उपरचना आहे की तिचा
  प्रांत $\Domain{N} = \Setabs{a\in \Domain{M}}{\Assign{R}{M}(a,
  b_1, \dots, b_n)}$ (पहा \ref{mod:bas:sub:rem:substructure}), आणि
  $\Expan{M}{X, b_1, \ldots, b_n}$ ही $R$, $c_1$, \dots, $c_n$ यांना
  अनुक्रमे $X$, $b_1,$ \dots, $b_n$ अशी अर्थनिर्धारणे देणारी
  $\Struct{M}$ ची विस्तारित रचना आहे (येथे $X \subseteq M^{n+1}$).
\end{enumerate}
\end{defn}

\begin{defn}
दोन अमूर्त तर्कशास्त्रे $\tuple{L_1, \models_{L_1}}$ आणि
$\tuple{L_2, \models_{L_2}}$ दिली असता, प्रत्येक भाषा
$\Lang{L}$ आणि वाक्य $\varphi_{E} \in L_1(\Lang{L})$ साठी असे
वाक्य $\varphi_{F} \in L_2(\Lang{L})$ असेल की $\Mod[L](L_1){\varphi_{E}} =
\Mod[L](L_2){\varphi_{F}}$, तर दुसरे तर्कशास्त्र पहिल्याइतके \emph{किमान
  तितकेच अभिव्यक्तिक्षम} आहे, असे म्हणतो आणि $\tuple{L_1, \models_{L_1}}
\leq \tuple{L_2, \models_{L_2}}$ असे लिहितो. $\tuple{L_1,
  \models_{L_1}} \leq \tuple{L_2, \models_{L_2}}$ आणि $\tuple{L_2,
  \models_{L_2}} \leq \tuple{L_1, \models_{L_1}}$ ही दोन्ही विधाने
खरी असतील, तर $\tuple{L_1, \models_{L_1}}$ आणि $\tuple{L_2,
  \models_{L_2}}$ ही तर्कशास्त्रे \emph{सममूल्य} आहेत.
\end{defn}

\begin{rem}
  प्रथम-क्रम तर्कशास्त्र, म्हणजे अमूर्त तर्कशास्त्र $\tuple{F,
  \models}$, हे सामान्य आहे. वस्तुतः, वरील गुणधर्म प्रथम-क्रम
  तर्कशास्त्रासाठी बहुतांशी सरळ आहेत. विस्तार गुणधर्म हा
  विस्तारात्मकतेवर येऊन ठेपतो, आणि वाक्य $\varphi_{E}$ चे
  $\Atom{R}{x, c_1, \dots, c_n}$ पर्यंतचे सापेक्षीकरण प्रत्येक
  उपसूत्र $\lforall[x][\varphi_{F}]$ च्या जागी
  $\lforall[x][(\Atom{R}{x, c_1, \dots, c_n} \lif \ \varphi_{F})]$ ठेवून
  मिळते, एवढेच आपण नोंदवतो. शिवाय, $\tuple{L, \models_L}$ सामान्य
  असेल, तर $\tuple{F, \models} \leq \tuple{L, \models_{L}}$; हे
  प्रथम-क्रम सूत्रे वरील विगमनाने दाखवता येते. म्हणून कोणतीही
  सामान्यता न गमावता प्रत्येक प्रथम-क्रम वाक्य प्रत्येक सामान्य
  तर्कशास्त्रात आहे असे आपण मानू शकतो.
\end{rem}












\section{संहतता आणि लोव्हेनहाइम--स्कोलेम गुणधर्म}\label{mod:lin:lsp:sec}

आता आपण संहतता आणि लोव्हेनहाइम--स्कोलेम यांचे अमूर्त तर्कशास्त्रांसाठीचे
स्वाभाविक विस्तार देऊ.

\begin{defn}
अमूर्त तर्कशास्त्र $\tuple{L, \models_L}$ ला \emph{संहतता गुणधर्म}
असतो, जर $L(\Lang{L})$-वाक्यांचा प्रत्येक संच $\Gamma$ असा
असेल की प्रत्येक सांत $\Gamma_0 \subseteq \Gamma$ पूर्ततायोग्य असताना
तो संच पूर्ततायोग्य असतो.
\end{defn}

\begin{defn}
प्रत्येक पूर्ततायोग्य $\Gamma$ ला गणनीय प्रतिमान असेल, तर
$\tuple{L, \models_L}$ ला \emph{अधोगामी लोव्हेनहाइम--स्कोलेम गुणधर्म}
असतो.
\end{defn}


\ref{mod:bas:pis:defn:partialisom} मधील आंशिक समरूपणाची संकल्पना निव्वळ
``बैजिक'' आहे (म्हणजे भाषेतील वाक्यांचा संदर्भ न घेता, केवळ
संरचनेच्या भाषा~$\Lang{L}$ ने दिलेल्या स्थिरांच्या
संदर्भात ती मांडलेली आहे), म्हणून ती अमूर्त तर्कशास्त्रांनाही लागू होते.
प्रथम-क्रम तर्कशास्त्राच्या बाबतीत, दोन संरचना आंशिकरीत्या समरूपी
असतील, तर त्या प्राथमिक सममूल्य असतात, हे आपल्याला
\ref{mod:bas:pis:thm:p-isom2} वरून माहीत आहे. मनमाना $\varphi_{E} \in
L(\Lang{L})$ साठी सूत्रे वरील विगमन उपलब्ध असेलच असे नसल्यामुळे ती
सिद्धता अमूर्त तर्कशास्त्रांना लागू होत नाही; तरीही लोव्हेनहाइम--स्कोलेम
गुणधर्म असेल, तर प्रमेय खरे असते.

\begin{thm}
\label{mod:lin:lsp:thm:abstract-p-isom}
$\tuple{L, \models_L}$ हे लोव्हेनहाइम--स्कोलेम गुणधर्म असलेले सामान्य
तर्कशास्त्र आहे असे समजा. मग कोणत्याही दोन आंशिकरीत्या समरूपी
संरचना या $\tuple{L, \models_L}$ मध्ये प्राथमिक सममूल्य असतात.
\end{thm}

\begin{proof}
$\Struct{M} \simeq_p \Struct{N}$ असे समजा; पण एखाद्या $\varphi_{E}$ साठी
$\Struct{M} \models_L \varphi_{E}$ आणि $\Struct{N} \not\models_L \varphi_{E}$ असेही
समजा. समरूपता गुणधर्मामुळे $\Domain{M}$ आणि $\Domain{N}$ हे परस्पर
वियुक्त आहेत असे आपण मानू शकतो; आणि विस्तार गुणधर्मामुळे सांत
भाषा~$\Lang{L}$ साठी $\varphi_{E} \in L(\Lang{L})$ आहे असे मानू शकतो.
$\Struct{M}$ आणि $\Struct{N}$ यांमधील आंशिक समरूपणांचा संच
$\mathcal{I}$ असू द्या; तसेच कोणतीही सामान्यता न गमावता, $p \in
\mathcal{I}$ आणि $q \subseteq p$ असेल, तर $q \in \mathcal{I}$ हेही
मानू.

$\Domain{M}^{<\omega}$ हा $\Domain{M}$ च्या घटक च्या सांत
क्रमिकांचा संच आहे. $S$ हा $\Domain{M}^{<\omega}$ वरील जोडणी दर्शवणारा
त्रिस्थानी संबंध असू द्या; म्हणजे $\mathbf{a}, \mathbf{b},
\mathbf{c} \in \Domain{M}^{<\omega}$ असतील, तर $S(\mathbf{a},
\mathbf{b}, \mathbf{c})$ तेव्हा आणि केवळ तेव्हाच खरे असते, जेव्हा
$\mathbf{c}$ ही $\mathbf{a}$ आणि $\mathbf{b}$ यांची जोडणी असते. तसेच
$T$ हा असा त्रिस्थानी संबंध असू द्या की $b \in M$ आणि
$\mathbf{a}, \mathbf{c} \in \Domain{M}^{<\omega}$ यांच्यासाठी
$T(\mathbf{a}, b, \mathbf{c})$ तेव्हा आणि केवळ तेव्हाच खरे असते,
जेव्हा $\mathbf{a} = a_1, \dots a_n$ आणि $\mathbf{c} = a_1, \dots
a_n, b$. नवीन 3-स्थानी विधेये $P$ आणि $Q$ निवडा आणि विश्व
$\Domain{M} \cup \Domain{M}^{<\omega}$ असलेली संरचना
$\Struct{M}^*$ बनवा; तिच्यात $\Struct{M}$ ही उपरचना असते आणि $P$ व
$Q$ ची अर्थनिर्धारणे अनुक्रमे जोडणी-संबंध $S$ व $T$ असतात (म्हणून
$\Struct{M}^*$ ही भाषा $\Lang{L} \cup \{ P, Q\}$ मधील आहे).

$\Domain{N}^{<\omega}$, $S'$, $T'$, $P'$, $Q'$ आणि $\Struct{N}^*$
यांची व्याख्या याच रीतीने करा. गृहीतकानुसार $\Struct{M} \simeq_p
\Struct{N}$ असल्यामुळे $\Domain{M}^{<\omega}$ आणि
$\Domain{N}^{<\omega}$ यांमध्ये असा संबंध $I$ असतो की
$I(\mathbf{a}, \mathbf{b})$ तेव्हा आणि केवळ तेव्हाच खरे असते, जेव्हा
$\mathbf{a}$ आणि $\mathbf{b}$ समरूपी असतात आणि
\ref{mod:bas:pis:defn:partialisom} मधील पुढे-मागे अट पूर्ण करतात. आता
$\Struct{A}$ ही अशी संरचना असू द्या की तिचा प्रांत हा
$\Struct{M}^*$ आणि $\Struct{N}^*$ यांच्या प्रांताचा संयोग आहे,
$\Struct{M}^*$ आणि $\Struct{N}^*$ या तिच्या उपसंरचना आहेत,
आणि तिच्या भाषा मध्ये संबंध~$I$ ने अर्थनिर्धारित केलेले एक
अतिरिक्त द्विपदी विधेय~$R$ व प्रांत $\Domain{M}^*$ आणि
$\Domain{N}^*$ दर्शवणारी विधेये आहेत.


\begin{figure}[h]
  \centering
  \begin{tikzpicture}[node distance=2cm, auto, thick, >=stealth']
    \draw [rounded corners] (0,0) -- (8,0) -- (8,4) -- (0,4) --  cycle;
    \draw (2,2) circle (0.5cm);
    \draw (2,2) circle (1.25cm);
    \draw (6,2) circle (0.5cm);
    \draw (6,2) circle (1.25cm);
    \path node at (0.75,3.5) {\large $\Struct{A}$};
    \path node at (2,2) {\large $\Struct{M}$};
    \path node at (6,2) {\large $\Struct{N}$};
    \path node at (3.5,1) {\large $\Struct{M}^*$};
    \path node at (7.5,1) {\large $\Struct{N}^*$};
    \node (Idom) at (2.8,2) {};
    \node (Irng) at (5.2,2) {};
    \draw[<->, bend left] (Idom) to node {\large $I$} (Irng) ;
  \end{tikzpicture}
  \caption{अंतर्गत आंशिक समरूपण असलेली संरचना~$\Struct{A}$.}
\end{figure}

निर्णायक निरीक्षण असे आहे की संरचना~$\Struct{A}$ च्या
भाषा मध्ये $\Struct{A}$ मध्ये सत्य असलेले असे
\emph{अमूर्त-तर्कशास्त्रीय} वाक्य $\theta_1$ असते की ते
$\Struct{M} \models_L \varphi_{E}$ आणि $\Struct{N} \not\models_L \varphi_{E}$ असे
सांगते (यासाठी सापेक्षीकरण गुणधर्म आवश्यक आहे); तसेच $\Struct{A}$ मध्ये
सत्य असलेले असे \emph{प्रथम-क्रम} वाक्य $\theta_2$ असते की ते
$I$ या आंशिक समरूपणाद्वारे $\Struct{M} \simeq_p \Struct{N}$ असे
सांगते. लोव्हेनहाइम--स्कोलेम गुणधर्मामुळे $\theta_1$ आणि $\theta_2$ ही दोन्ही
एका गणनीय प्रतिमान $\Struct{A}_0$ मध्ये सत्य असतात; त्यात
आंशिकरीत्या समरूपी अशा उपरचना $\Struct{M}_0$ आणि $\Struct{N}_0$ असतात,
ज्यांच्यासाठी $\Struct{M}_0 \models_L \varphi_{E}$ आणि $\Struct{N}_0
\not\models_L \varphi_{E}$. पण गणनीय आणि आंशिकरीत्या समरूपी
संरचना या \ref{mod:bas:pis:thm:p-isom1} नुसार वस्तुतः समरूपी
असतात; हे सामान्य अमूर्त तर्कशास्त्रांच्या समरूपता गुणधर्माच्या विरुद्ध
आहे.


\end{proof}












\section{लिंडस्ट्रॉमचे प्रमेय}\label{mod:lin:prf:sec}

\begin{lem}
\label{mod:lin:prf:lem:lindstrom}
$\varphi_{E} \in L(\Lang{L})$ आणि $\Lang{L}$ सांत आहे असे समजा; तसेच असा
$n \in \Nat$ आहे असे समजा की कोणत्याही दोन संरचना
$\Struct{M}$ आणि~$\Struct{N}$ साठी, $\Struct{M} \equiv_n
\Struct{N}$ आणि $\Struct{M} \models_L \varphi_{E}$ असेल, तर $\Struct{N}
\models_L \varphi_{E}$ हेही खरे असते. मग $\varphi_{E}$ हे एका प्रथम-क्रम
वाक्य शी सममूल्य असते; म्हणजे असे प्रथम-क्रम $\theta$ असते की
$\Mod(L){\varphi_{E}} = \Mod(L){\theta}$.
\end{lem}

\begin{proof}
कोणत्याही दोन $n$-सममूल्य संरचना $\Struct{M}$ आणि
$\Struct{N}$ या $\varphi_{E}$ ला नेमलेल्या सत्यतामूल्याबाबत सहमत असतील, असा $n$
घ्या. \ref{mod:bas:pis:prop:qr-finite} आठवा: तार्किक सममूल्यतेपर्यंत,
सांत भाषा मध्ये संख्यापक-प्रतांक $n$ पेक्षा अधिक नसलेली प्रथम-क्रम
वाक्ये सांत इतकीच असतात. आता प्रत्येक स्थिर
संरचना~$\Struct{M}$ साठी, $\Struct{M}$ मध्ये सत्य आणि
$\QuantRank{\varphi_{E}} \le n$ असलेल्या सर्व प्रथम-क्रम वाक्ये~$\varphi_{E}$
यांची संधी $\theta_{\Struct{M}}$ असू द्या (ही संधी सांत आहे), म्हणजे
$\Struct{N} \models \theta_{\Struct{M}}$ तेव्हा आणि केवळ तेव्हाच, जेव्हा
$\Struct{N} \equiv_n \Struct{M}$. मग $\theta = \textstyle\bigvee
\Setabs{\theta_{\Struct{M}}}{\Struct{M} \models_L \varphi_{E}}$ असे ठेवा; हा विकल्पही
(तार्किक सममूल्यतेपर्यंत) सांत आहे.

यावरून $\Mod(L){\varphi_{E}} = \Mod(L){\theta}$ हा निष्कर्ष मिळतो. वस्तुतः,
$\Struct{N} \models_L \theta$ असेल, तर एखाद्या $\Struct{M} \models_L
\varphi_{E}$ साठी $\Struct{N} \models \theta_{\Struct{M}}$ असते; त्यामुळे
लेम्माच्या गृहीतकानुसार $\Struct{N} \models_L \varphi_{E}$ हेही खरे असते.
उलट, $\Struct{N} \models_L \varphi_{E}$ असेल, तर $\theta_\Struct{N}$ हा $\theta$ मधील
एक विकल्पांश आहे; आणि $\Struct{N} \models \theta_\Struct{N}$ असल्यामुळे
$\Struct{N} \models_L \theta$ हेही खरे असते.
\end{proof}

\begin{thm}[लिंडस्ट्रॉमचे प्रमेय]
  \label{mod:lin:prf:thm:lindstrom} सामान्य अमूर्त तर्कशास्त्र
  $\tuple{L, \models_L}$ ला संहतता आणि लोव्हेनहाइम--स्कोलेम गुणधर्म
  आहेत असे समजा. मग $\tuple{L, \models_L} \le \tuple{F, \models}$
  (म्हणून $\tuple{L, \models_L}$ हे प्रथम-क्रम तर्कशास्त्राशी सममूल्य
  आहे).

\end{thm}

\begin{proof}
\ref{mod:lin:prf:lem:lindstrom} नुसार, सांत $\Lang{L}$ साठी कोणतेही $\varphi_{E}
\in L(\Lang{L})$ दिले असता असा $n \in \Nat$ असतो की कोणत्याही दोन
संरचना $\Struct{M}$ आणि~$\Struct{N}$ साठी: $\Struct{M}
\equiv_n \Struct{N}$ असेल, तर $\Struct{M}$ आणि $\Struct{N}$ या $\varphi_{E}$
बाबत सहमत असतात, हे दाखवणे पुरेसे आहे. कारण मग $\varphi_{E}$ हे प्रथम-क्रम
वाक्य शी सममूल्य असते, आणि त्यावरून $\tuple{L, \models_L} \le
\tuple{F, \models}$ मिळते. आपण सांत आणि निव्वळ संबंधात्मक
भाषा मध्ये काम करत असल्यामुळे, \ref{mod:bas:pis:thm:b-n-f}
नुसार $\Struct{M} \equiv_n \Struct{N}$ या विधानाच्या जागी अनुरूप बैजिक
विधान $I_n(\emptyset,\emptyset)$ घेता येते.

$\varphi_{E}$ दिले असता, विरोधासाठी असे समजा की प्रत्येक $n$ साठी अशा
संरचना $\Struct{M}_n$ आणि $\Struct{N}_n$ असतात की
$I_n(\emptyset, \emptyset)$; पण (समजा) $\Struct{M}_n \models_L \varphi_{E}$
आणि $\Struct{N}_n \not\models_L \varphi_{E}$. भाषेत सांत इतकीच आण्विक
वाक्ये असल्यामुळे, गरज पडल्यास $\Struct{M}_n$ चा उपअनुक्रम
घेऊन त्या सर्व त्याच आण्विक वाक्यांची पूर्ति करतात असे मानू
शकतो. हा उपअनुक्रम अनंत निवडून नव्याने क्रमांकित केल्यास प्रत्येक मूळ
निर्देशांक नव्या निर्देशांकाएवढा किंवा त्याहून मोठा असतो; म्हणून आवश्यक
सममूल्यतेची खोली कायम राहते. आता स्थिरांतील समानतेचा आकृतिबंध सर्वत्र
सारखा आहे. त्यामुळे समरूपता गुणधर्म वापरून सर्व $\Struct{M}_n$ भाषेतील
स्थिरांना त्याच वस्तू नेमतात आणि त्या वस्तूंखेरीज त्यांचे प्रांत परस्पर
विभक्त आहेत असे मानू शकतो. सर्व $\Struct{M}_n$ चा संयोग,
म्हणजे प्रत्येक $\Struct{M}_n$ उपरचना असलेली एकमेव न्यूनतम
संरचना, $\Struct{M}$ असू द्या. \ref{mod:lin:lsp:thm:abstract-p-isom}
च्या सिद्धतेप्रमाणे, $\Struct{M}^*$ ही $\Struct{M}$ ची अशी वर्धित रचना
असू द्या की तिचा प्रांत $\Domain{M} \cup
\Domain{M}^{<\omega}$ आहे आणि तिच्या विस्तारित भाषा मध्ये
जोडणी-विधेये $P$ आणि~$Q$ आहेत.




$\Struct{N}_n$, $\Struct{N}$ आणि $\Struct{N}^*$ यांची व्याख्या याच
रीतीने करा. आता $\Struct{A}$ ही अशी संरचना असू द्या की तिचा
प्रांत हा $\Struct{M}^*$ आणि $\Struct{N}^*$ यांच्या प्रांत
बरोबर नैसर्गिक संख्या~$\Nat$ आणि त्यांचा नैसर्गिक क्रम~$\le$ यांचा
समावेश करतो. तिच्या भाषा मध्ये प्रांत $\Domain{M}$,
$\Domain{N}$, $\Domain{M}^{<\omega}$ आणि $\Domain{N}^{<\omega}$
दर्शवणारी अतिरिक्त विधेये असतात; तसेच $\Struct{M}_n$ आणि
$\Struct{N}_n$ यांचे प्रांत पुढील अर्थाने संकेतबद्ध करणारी विधेये असतात:
\begin{align*}
  \Domain{M_n} & = \Setabs{a \in \Domain{M}}{R(a, n)}; &
  \Domain{N_n} & = \Setabs{a \in \Domain{N}}{S(a,n)}; \\
  \Domain{M}^{<\omega}_n & = \Setabs{a \in \Domain{M}^{<\omega}}{R(a,n)}; &
  \Domain{N}^{<\omega}_n & = \Setabs{a \in \Domain{N}^{<\omega}}{S(a,n)}.
\end{align*}
संरचना~$\Struct{A}$ मध्ये असा त्रिस्थानी संबंध $J$ ही असतो की
$J(n, \mathbf{a}, \mathbf{b})$ तेव्हा आणि केवळ तेव्हाच खरे असते,
जेव्हा $I_n(\mathbf{a}, \mathbf{b})$.


आता $R$, $S$, $J$ इत्यादींनी वाढवलेल्या भाषा~$\Lang{L}$ मध्ये
असे वाक्य~$\theta$ असते की ते $\le$ हा पहिला घटक असलेला पण शेवटचा
घटक नसलेला विविक्त रेषीय क्रम आहे, $\Struct{M}_n \models_L \varphi_{E}$,
$\Struct{N}_n \not\models_L \varphi_{E}$, आणि क्रमातील प्रत्येक $n$ साठी
$J(n, \mathbf{a}, \mathbf{b})$ तेव्हा आणि केवळ तेव्हाच खरे असते,
जेव्हा $I_n(\mathbf{a}, \mathbf{b})$, असे सांगते.


संहतता गुणधर्म वापरून, $\theta$ बरोबर एका नवीन स्थिराला प्रत्येक मानक
घटकाच्या वर ठेवणारी वाक्ये घ्या. या संचाला असे प्रतिमान $\Struct{A}^*$
मिळते की त्यातील क्रमात अमानक घटक~$n^*$ असतो. विशेषतः,
$\Struct{A}^*$ मध्ये अशा उपसंरचना $\Struct{M_{n^*}}$ आणि
$\Struct{N_{n^*}}$ असतात की $\Struct{M_{n^*}} \models_L \varphi_{E}$ आणि
$\Struct{N_{n^*}} \not\models_L \varphi_{E}$. आता $\Domain{M_{n^*}}$ मधील
आणि $\Domain{N_{n^*}}$ मधील $k$-सदस्यी क्रमितकांच्या जोड्यांचा संच
$\mathcal{I}$ पुढीलप्रमाणे परिभाषित करता येतो: $\tuple{\mathbf{a},
\mathbf{b}} \in \mathcal{I}$ तेव्हा आणि केवळ तेव्हाच, जेव्हा
$J(n^*-k, \mathbf{a}, \mathbf{b})$; येथे $k$ ही $\mathbf{a}$ आणि
$\mathbf{b}$ यांची लांबी आहे. $n^*$ हा अमानक असल्यामुळे प्रत्येक मानक
$k$ साठी $n^* - k >0$ असते, आणि $\mathcal{I}$ हा संच
$\Struct{M_{n^*}} \simeq_p \Struct{N_{n^*}}$ या वस्तुस्थितीचा साक्षी
आहे. पण \ref{mod:lin:lsp:thm:abstract-p-isom} नुसार $\Struct{M_{n^*}}$ ही
$\Struct{N_{n^*}}$ शी $L$-सममूल्य आहे; हा विरोधाभास आहे.


\end{proof}



